%=========== ΛΥΣΗ - 26ο ΘΕΜΑ Α ===========
\providecommand{\theoriaref}[3]{%
\begin{center}
\fcolorbox{maincolor!70!black}{maincolor!8}{%
\begin{minipage}{.88\linewidth}
\textcolor{maincolor!80!black}{\kerkissans{\textbf{#1~\ref{#2}}}}\\[1mm]
#3
\end{minipage}}
\end{center}}
\providecommand{\orismosref}[1]{%
\theoriaref{Ορισμός}{#1}{Η απάντηση δίνεται από τον αντίστοιχο ορισμό.}}
\providecommand{\apodeixiref}[1]{%
\theoriaref{Θεώρημα}{#1}{Η ζητούμενη απόδειξη δίνεται στην αντίστοιχη απόδειξη.}}

\begin{thema}{Α}
\begin{erwthma}
\item \apodeixiref{thewrhma:kritirio_topikwn_akrotatwn}

\item \orismosref{orismos:gnisiws_fthinousa}

\item \orismosref{orismos:orizontia_asymptwth}

\item Οι χαρακτηρισμοί των προτάσεων είναι:
\begin{alist}
\item \textbf{Λάθος}. Για να υπάρχει το $\lim\limits_{x\to x_0}f(x)=\ell$, δεν είναι απαραίτητο η $f$ να ορίζεται στο $x_0$.
\item \textbf{Λάθος}. Κρίσιμα σημεία είναι τα εσωτερικά σημεία όπου μηδενίζεται η $f'$ ή όπου δεν ορίζεται η $f'$.
\item \textbf{Σωστό}. Ισχύει
\[ (\ln{x})'=\dfrac{1}{x},\quad x>0 \]
και
\[ (\ln{|x|})'=\dfrac{1}{x},\quad x\neq0. \]
\item \textbf{Σωστό}. Αν η $f$ είναι γνησίως αύξουσα και αντιστρέφεται, τότε και η $f^{-1}$ είναι γνησίως αύξουσα.
\item \textbf{Λάθος}. Θετικό ολοκλήρωμα δεν σημαίνει ότι η $f$ είναι θετική σε κάθε σημείο του $[a,\beta]$. Μπορεί να παίρνει και αρνητικές τιμές, αλλά το συνολικό προσημασμένο εμβαδόν να είναι θετικό.
\end{alist}

\item Σωστή απάντηση είναι η \textbf{γ}. Για μια συνάρτηση $1-1$ ισχύει
\[ f(x_1)=f(x_2)\Rightarrow x_1=x_2. \]
\end{erwthma}
\end{thema}
%=========== ΤΕΛΟΣ ΛΥΣΗΣ - 26ο ΘΕΜΑ Α ===========
